% ===========================================================================
%  9 — RAG e camada de linguagem
% ===========================================================================
\chapter{Recuperação documental, geração e agente}
\label{cap:rag}

Este capítulo documenta a camada que transforma um diagnóstico em
\emph{instrução de correção}. É onde o requisito mais exigente do enunciado
--- responder apenas sobre falhas documentadas --- é tecnicamente garantido.

\section{Visão do pipeline}

\begin{figure}[H]
\centering
\fitwidth{%
\begin{tikzpicture}[x=1cm, y=1cm]

  % ---- trilha de indexação (offline)
  \node[c4boundarylbl, anchor=north west, fill=pmxSurface] at (0.3,0.35)
    {INDEXAÇÃO (offline ou no primeiro \textit{boot})};
  \draw[c4boundary, fill=pmxBlue!4] (0.2,0.0) rectangle (17.0,-3.4);

  \node[etapaio, text width=22mm] (pdf) at (1.8,-1.7)
    {\arq{DocN.pdf}\\{\scriptsize procedimento}};
  \node[etapa, text width=24mm] (ext) at (5.2,-1.7)
    {\textbf{Extração}\\{\scriptsize página a página; rastreia seção corrente}};
  \node[etapaerr, text width=22mm] (ocr) at (8.5,-1.7)
    {\textbf{OCR}\\{\scriptsize só se não houver camada de texto}};
  \node[etapa, text width=24mm] (chk) at (11.8,-1.7)
    {\textbf{\textit{Chunking}}\\{\scriptsize $\approx$\num{1200} car.,
     sobrep.\ 150}};
  \node[etapa, text width=24mm] (emb) at (15.1,-1.7)
    {\textbf{\textit{Embedding}}\\{\scriptsize dim.\ 768}};

  \draw[pipe] (pdf) -- (ext);
  \draw[pipe] (ext) -- (ocr) node[pipelbl, midway, above=1pt]{sem texto};
  \draw[pipe] (ocr) -- (chk);
  \draw[pipe] (ext.north) -- (5.2,-0.8) -- (11.8,-0.8) -- (chk.north)
    node[pipelbl, pos=0.5, above=1pt]{com camada de texto};
  \draw[pipe] (chk) -- (emb);

  % ---- base vetorial
  \node[c4db, text width=30mm, minimum height=15mm] (ch) at (8.6,-5.6)
    {\textbf{ChromaDB --- coleção \cd{manuals}}\tech{HNSW, cosseno}};
  \draw[pipe] (emb.south) -- (15.1,-5.6) -- (ch.east)
    node[pipelbl, pos=0.7, above=1pt]{um conjunto por família coberta};

  % ---- trilha de consulta (online)
  \node[c4boundarylbl, anchor=north west, fill=pmxSurface] at (0.3,-7.15)
    {CONSULTA (tempo de requisição)};
  \draw[c4boundary, fill=pmxGreen!4] (0.2,-7.5) rectangle (17.0,-13.6);

  \node[etapaio, text width=24mm] (fam) at (2.0,-8.9)
    {\textbf{Família prevista}\\{\scriptsize saída do motor de diagnóstico}};
  \node[decision, text width=30mm] (gd) at (6.4,-8.9)
    {\textbf{Possui \textit{chunk} indexado?}};
  \node[etapaerr, text width=32mm] (nao) at (12.4,-8.9)
    {\textbf{Aviso de falha não documentada}\\{\scriptsize + sugestão de
     registro --- \textbf{o modelo não é chamado}}};
  \node[etapa, text width=26mm] (ret) at (4.2,-12.2)
    {\textbf{Recuperação}\\{\scriptsize \textit{top-k}=6, filtro por família}};
  \node[etapa, text width=26mm] (pr) at (8.6,-12.2)
    {\textbf{\textit{Prompt} de contexto fechado}\\{\scriptsize trechos
     numerados}};
  \node[etapaout, text width=28mm] (ger) at (13.2,-12.2)
    {\textbf{Instruções + citações}\\{\scriptsize citação dos metadados, não do
     texto gerado}};

  \draw[pipe] (fam) -- (gd);
  \draw[pipe, draw=pmxOrange] (gd.east) -- (nao.west)
    node[pipelbl, midway, above=1pt]{não};
  \draw[pipe] (gd.south) -- (6.4,-10.6) -- (4.2,-10.6) -- (ret.north)
    node[pipelbl, pos=0.25, above=1pt]{sim};
  \draw[pipe] (ret) -- (pr);
  \draw[pipe] (pr) -- (ger);
  \draw[c4reldash] (ch.south) -- (8.6,-6.6) -- (4.2,-6.6) -- (4.2,-8.2)
    -- (ret.north);
  \draw[c4reldash] (ch.west) -- (6.4,-5.6) -- (6.4,-7.9)
    node[cxlabel, pos=0.85, right=2pt]{consulta de cobertura};

\end{tikzpicture}}
\caption[Pipeline de recuperação e geração]{Pipeline completo de recuperação e
  geração. As duas trilhas compartilham a mesma base vetorial, que cumpre dois
  papéis: repositório dos trechos \emph{e} fonte de verdade da cobertura
  documental. O ramo laranja é o caminho exigido pelo enunciado --- e nele o
  modelo de linguagem não é invocado.}
\label{fig:ragpipe}
\end{figure}

\section{Ingestão documental}

\subsection{Extração e \textit{chunking}}

A extração percorre o PDF página a página, mantendo o rastreio da seção
corrente por reconhecimento de títulos numerados no padrão \cd{N. Título} ou
\cd{N.N Título}. Os parágrafos são acumulados em \textit{chunks} de
aproximadamente \num{1200} caracteres com sobreposição de 150.

\begin{tabularx}{\linewidth}{@{}L{34mm}L{20mm}X@{}}
\toprule
\headrow \thd{Parâmetro} & \thd{Valor} & \thd{Justificativa}\\
\midrule
Tamanho do \textit{chunk} & $\approx$\num{1200} caracteres &
  Grande o bastante para conter um passo de procedimento completo com seu
  contexto; pequeno o bastante para que o \textit{top-k} de 6 trechos caiba
  confortavelmente no \textit{prompt}.\\
\zebra Sobreposição & 150 caracteres &
  Evita que um passo de procedimento seja cortado exatamente na fronteira entre
  dois \textit{chunks}, o que o tornaria irrecuperável.\\
Descarte mínimo & 80 caracteres &
  Fragmentos menores são cabeçalhos, números de página e rodapés --- ruído que
  degradaria a busca.\\
\zebra Metadados & doc, título, seção, página, família &
  \textbf{São a citação.} O usuário vê estes campos, não uma referência
  produzida pelo modelo.\\
\bottomrule
\end{tabularx}

\subsection{OCR: o caso do documento escaneado}

Um dos procedimentos da base inicial é um PDF escaneado, sem camada de texto ---
situação comum em acervos industriais, onde procedimentos antigos existem apenas
em papel digitalizado. A extração convencional devolve zero \textit{chunks}.

O tratamento é um \textit{fallback} explícito: se a extração de texto não
produzir \textit{chunk} algum \emph{e} houver um transcritor disponível, cada
página é renderizada em imagem a \SI{150}{dpi} e transcrita por modelo
multimodal, com instrução para preservar títulos numerados em linhas próprias
--- o que mantém o rastreio de seção funcionando sobre o texto transcrito.

\begin{decisao}
O OCR é \emph{condicional}, não incondicional. Aplicá-lo sempre custaria uma
chamada de modelo por página em todo documento --- caro e mais sujeito a erro
que a extração direta, que é exata quando a camada de texto existe. A
degradação só ocorre quando é necessária.
\end{decisao}

\section{\textit{Embeddings} e base vetorial}

\begin{tabularx}{\linewidth}{@{}L{34mm}X@{}}
\toprule
\headrow \thd{Aspecto} & \thd{Decisão e justificativa}\\
\midrule
Modelo & \cd{gemini-embedding-2}, multilíngue, contexto de \num{8192}
  \textit{tokens}.\\
\zebra Dimensionalidade & 768, obtida por truncamento de representação
  aninhada~\cite{kusupati2022matryoshka}. A propriedade dessas representações é
  que os primeiros componentes concentram a informação mais relevante, de modo
  que truncar não exige retreino e degrada pouco a qualidade --- em troca de
  um quarto do armazenamento e da latência de busca.\\
Tipo de tarefa & \cd{RETRIEVAL\_DOCUMENT} na indexação e
  \cd{RETRIEVAL\_QUERY} na consulta. A assimetria é intencional: documento e
  pergunta são objetos linguísticos diferentes, e o modelo os projeta em
  espaços alinhados mas com condicionamento distinto.\\
\zebra Métrica e índice & Cosseno sobre índice HNSW~\cite{malkov2020hnsw}, que
  oferece busca aproximada com complexidade logarítmica --- irrelevante nos
  268 \textit{chunks} atuais, decisivo em escala de milhares de
  procedimentos.\\
Lote de requisição & Uma instância por requisição. O serviço aceita lista, mas
  retorna silenciosamente apenas um vetor; o cliente valida explicitamente que
  o número de vetores devolvidos é igual ao de textos enviados e levanta
  exceção caso contrário.\\
\bottomrule
\end{tabularx}

\begin{atencao}
A validação de contagem de \textit{embeddings} parece defensiva em excesso até
se considerar a falha que ela previne: um desalinhamento silencioso entre
textos e vetores associaria o \textit{embedding} de um trecho ao conteúdo de
outro. O resultado seria um sistema que recupera trechos errados com total
confiança --- e sem nenhuma mensagem de erro.
\end{atencao}

\section{O \textit{guardrail} de cobertura}

\subsection{Definição operacional}

\begin{especificacao}{Regra de cobertura documental}
Uma família de falha é considerada \textbf{documentada} se, e somente se,
existe ao menos um \textit{chunk} indexado na base vetorial cujo metadado
\cd{family} é igual ao nome canônico dessa família.
\end{especificacao}

Trata-se de uma consulta ao índice, com filtro e limite 1 --- não de uma
inferência. Três consequências decorrem dessa formulação:

\begin{enumerate}
  \item \textbf{É verificável.} O estado da cobertura é inspecionável a qualquer
    momento pelos \textit{endpoints} de saúde e de documentos.
  \item \textbf{É dinâmica.} Cadastrar um procedimento altera a cobertura na
    requisição seguinte, sem reinício de serviço e sem edição de configuração.
  \item \textbf{Não pode ser contornada pelo modelo.} A verificação ocorre no
    fluxo de controle, antes da montagem do \textit{prompt}. Não existe caminho
    de código em que o modelo seja chamado para uma família descoberta.
\end{enumerate}

\subsection{Cobertura inicial e as lacunas deliberadas}

\begin{table}[H]
\caption{Cobertura documental inicial}
\label{tab:cobertura}
\begin{tabularx}{\linewidth}{@{}L{22mm}L{52mm}X@{}}
\toprule
\headrow \thd{Documento} & \thd{Famílias cobertas} & \thd{Tema}\\
\midrule
\arq{Doc1.pdf} & \cd{rolamento\_inner}, \cd{rolamento\_outer},
  \cd{rolamento\_ball}, \cd{rolamento\_combination} &
  Diagnóstico e correção em rolamentos (documento escaneado --- exige OCR)\\
\zebra \arq{Doc2.pdf} & \cd{desalinhamento} & Correção de desalinhamento em
  motor elétrico\\
\arq{Doc3.pdf} & \cd{desbalanceamento} & Correção de desbalanceamento em
  máquinas rotativas\\
\zebra \arq{Doc4.pdf} & \cd{correia} & Transmissão por correias\\
\arq{Doc5.pdf} & \cd{polia} & Polias de sistemas rotativos\\
\zebra \arq{Doc6.pdf} & \cd{cocked\_rotor} & Rotor inclinado\\
\midrule
\multicolumn{3}{@{}l}{\textbf{Sem documento --- acionam o fluxo de falha não
  documentada}}\\
--- & \cd{eccentric\_rotor} & \num{16497} ocorrências --- maior lacuna
  documental do acervo\\
\zebra --- & \cd{ventoinha} & \num{12299} ocorrências\\
--- & \cd{falta\_fase} & 800 ocorrências --- e a classe que o modelo melhor
  reconhece\\
\bottomrule
\end{tabularx}
\end{table}

Nove das doze famílias de falha estão cobertas por 268 \textit{chunks}. As três
lacunas não são acidentais: elas permitem demonstrar o comportamento exigido
pelo enunciado com dados reais, e alimentam a matriz de priorização documental
do painel (Seção~\ref{sec:priorizacao}).

\section{Anti-alucinação: cinco barreiras}

A geração de conteúdo plausível mas não sustentado pela fonte é a falha
característica de modelos de linguagem~\cite{ji2023hallucination}. O padrão de
recuperação aumentada~\cite{lewis2020rag} ataca a raiz do problema ---
substituir conhecimento paramétrico por conhecimento recuperado --- mas não é
suficiente por si só. O sistema aplica cinco barreiras em série.

\begin{table}[H]
\caption{Barreiras anti-alucinação}
\label{tab:antialuc}
\begin{tabularx}{\linewidth}{@{}L{8mm}L{40mm}X@{}}
\toprule
\headrow \thd{\#} & \thd{Barreira} & \thd{Mecanismo}\\
\midrule
1 & \textit{Guardrail} determinístico & O modelo só é chamado se a família
  possui trechos indexados. A cobertura nunca é decidida pelo modelo.\\
\zebra 2 & Contexto fechado & A instrução de sistema exige responder
  \textbf{somente} com base nos trechos fornecidos e declarar explicitamente
  quando eles não cobrem o ponto perguntado.\\
3 & Citação determinística & Documento, seção e página vêm dos metadados do
  trecho recuperado, não do texto gerado. \textbf{Citação alucinada é
  estruturalmente impossível.}\\
\zebra 4 & Temperatura baixa & \cd{temperature=0.2} na geração;
  \cd{0.0} na transcrição por OCR, onde qualquer criatividade é defeito.\\
5 & Confiança dupla exposta & Probabilidade $\times$ concordância dos vizinhos
  é apresentada ao usuário; abaixo do limiar, a interface marca o diagnóstico
  como incerto.\\
\bottomrule
\end{tabularx}
\end{table}

\subsection{A instrução de sistema}

\begin{lstlisting}[language=Python, caption={Instrução de sistema do assistente prescritivo --- \arq{src/llm/prompts.py}}, label={lst:sysprompt}]
SYSTEM_PRESCRIPTIVE = """Voce e um assistente tecnico de manutencao industrial de maquinas rotativas.
Regras OBRIGATORIAS:
1. Responda SOMENTE com base nos trechos de documentos fornecidos no contexto.
2. Se a informacao nao estiver no contexto, diga explicitamente que os documentos
   disponiveis nao cobrem esse ponto.
3. Ao usar um trecho, referencie-o pelo marcador [n] correspondente.
4. Nunca invente procedimentos, valores de tolerancia, normas ou frequencias que
   nao estejam no contexto.
5. Responda em portugues do Brasil, em Markdown, com passos numerados quando
   descrever procedimentos.
6. Priorize sempre a secao de seguranca antes de qualquer intervencao."""
\end{lstlisting}

A sexta regra não é anti-alucinação --- é segurança operacional. Um
procedimento de manutenção que descreve a intervenção antes do bloqueio e
etiquetagem do equipamento está tecnicamente correto e operacionalmente
perigoso.

\subsection{Estrutura obrigatória da prescrição}

O \textit{prompt} de prescrição exige cinco seções, na ordem em que um técnico
as executa: \textbf{segurança} (passos obrigatórios antes da intervenção),
\textbf{diagnóstico} (como confirmar a falha, incluindo as assinaturas de
vibração citadas nos trechos), \textbf{correção} (procedimento passo a passo),
\textbf{validação} (como verificar e quais são os critérios de aceitação) e
\textbf{prevenção}. A estrutura fixa torna a saída comparável entre falhas e
auditável.

\section{Agente de conversa}

\subsection{Fronteira arquitetural}

O agente existe \textbf{apenas} no \textit{endpoint} de conversa
(\adrr{ADR-11}). O caminho crítico de diagnóstico --- usado pelo
\textit{worker} MQTT, e portanto por toda a planta --- permanece
determinístico. A razão é direta: um laço agêntico tem número de passos
variável, latência imprevisível e liberdade de decisão que enfraqueceria
exatamente as garantias descritas na seção anterior. Onde há garantia a
preservar, não há agente.

\subsection{Ferramentas}

\begin{figure}[H]
\centering
\fitwidth{%
\begin{tikzpicture}[x=1cm, y=1cm]

  \node[c4container, text width=30mm, minimum height=13mm] (ui) at (2.0,0)
    {\textbf{Conversa}\tech{Streamlit}\desc{sessão por aba}};
  \node[c4container, text width=30mm, minimum height=13mm] (ep) at (6.6,0)
    {\textbf{\cd{POST /chat}}\tech{FastAPI}};
  \node[c4component, text width=36mm, minimum height=15mm] (ag) at (11.8,0)
    {\textbf{Agente}\tech{LlmAgent + Runner}\desc{memória de sessão no
     Postgres, esquema \cd{adk}}};
  \node[c4component, text width=32mm, minimum height=13mm,
        fill=pmxYellowWash, draw=pmxYellow!70!black] (fb) at (11.8,-3.4)
    {\textbf{\textit{Fallback}: RAG de um passo}\desc{sem agente ou sem
     credenciais}};

  \node[c4component, text width=40mm, minimum height=11mm] (t1) at (18.4,1.9)
    {\textbf{\cd{consultar\_historico}}\desc{SQL no Postgres}};
  \node[c4component, text width=40mm, minimum height=11mm,
        fill=pmxOrangeWash, draw=pmxOrange] (t2) at (18.4,0.2)
    {\textbf{\cd{buscar\_documentos}}\desc{\textit{guardrail} + recuperação}};
  \node[c4component, text width=40mm, minimum height=11mm] (t3) at (18.4,-1.5)
    {\textbf{\cd{diagnosticar\_evento}}\desc{motor de diagnóstico}};
  \node[c4component, text width=40mm, minimum height=11mm] (t4) at (18.4,-3.2)
    {\textbf{\cd{cobertura\_documental}}\desc{índice + banco}};

  \draw[c4rel] (ui) -- (ep);
  \draw[c4rel] (ep) -- (ag) node[cxlabel, midway, above=1pt]{agente disponível};
  \draw[c4rel] (ep.south) -- (6.6,-3.4) -- (fb.west)
    node[cxlabel, pos=0.62, above=1pt]{degradação automática};
  \draw[c4rel] (ag.east) -- (16.3,1.9);
  \draw[c4rel] (ag.east) -- (16.3,0.2);
  \draw[c4rel] (ag.east) -- (16.3,-1.5);
  \draw[c4rel] (ag.east) -- (16.3,-3.2);

\end{tikzpicture}}
\caption[Agente de conversa e ferramentas]{Agente de conversa e suas quatro
  ferramentas. Em laranja, a ferramenta que carrega o \textit{guardrail}: a
  verificação de cobertura acontece \emph{dentro} dela, antes da recuperação, de
  modo que o modelo recebe um aviso estruturado em vez de trechos --- e apenas o
  transmite.}
\label{fig:agente}
\end{figure}

\begin{table}[H]
\caption{Ferramentas do agente}
\label{tab:ferramentas}
\begin{tabularx}{\linewidth}{@{}L{38mm}X@{}}
\toprule
\headrow \thd{Ferramenta} & \thd{Contrato e comportamento}\\
\midrule
\cd{consultar\_historico(familia, dias)} &
  Consulta agregada ao histórico. Valida a família contra a taxonomia e devolve
  a lista de famílias válidas em caso de erro --- o que permite ao modelo se
  autocorrigir. A janela temporal é ancorada no evento \emph{mais recente do
  banco}, não na data corrente, o que mantém a resposta coerente em um conjunto
  de dados histórico.\\
\zebra \cd{buscar\_documentos(consulta, familia)} &
  \textbf{Obrigatória} antes de descrever qualquer procedimento. Aplica o
  \textit{guardrail}: família sem cobertura devolve
  \cd{documentado=false} e o aviso pronto. Família coberta devolve os trechos
  numerados com origem. Registra as citações no coletor de contexto.\\
\cd{diagnosticar\_evento(json)} &
  Executa o motor de diagnóstico sobre um JSON colado pelo usuário. JSON
  inválido devolve os erros de validação em forma legível, não uma exceção.\\
\zebra \cd{cobertura\_documental()} &
  Lista famílias documentadas, famílias sem documento e documentos cadastrados
  com data de ingestão.\\
\bottomrule
\end{tabularx}
\end{table}

\subsection{Citações determinísticas no agente}

Em um agente, a dificuldade adicional é que a recuperação acontece dentro de uma
ferramenta, várias camadas abaixo do código que monta a resposta HTTP. A solução
adotada é um coletor por contexto de execução (\textit{context variable}): a
ferramenta registra ali os metadados dos trechos que efetivamente recuperou, e o
\textit{endpoint} lê o coletor ao final da execução.

\begin{decisao}
O coletor por contexto de execução é preferível a uma variável global porque é
isolado por tarefa assíncrona: duas conversas simultâneas não misturam citações.
E é preferível a inspecionar o texto gerado pelo modelo em busca de marcadores,
porque preserva a propriedade central --- a citação exibida reflete o que foi
\emph{recuperado}, não o que o modelo \emph{escreveu}.
\end{decisao}

\subsection{Degradação graciosa}

A inicialização do agente pode falhar por três motivos: biblioteca ausente,
projeto de nuvem não configurado ou credenciais inválidas. Em todos os casos a
função de obtenção do agente devolve nulo, e o \textit{endpoint} de conversa
recorre ao RAG de um passo. A interface indica o modo em uso por um selo
--- \enquote{agente com ferramentas e memória} ou \enquote{recuperação direta}
--- para que o usuário saiba o que esperar.

\begin{figure}[H]
\centering
\fitwidth{%
\begin{tikzpicture}[x=1cm, y=1cm]

  \raia{1.7}{-9.6}{Usuário}{interface}
  \raia{5.4}{-9.6}{\cd{POST /chat}}{FastAPI}
  \raia{9.1}{-9.6}{Agente}{LlmAgent}
  \raia{12.8}{-9.6}{Ferramentas}{4 funções}
  \raiaext{16.5}{-9.6}{Recursos}{Postgres, Chroma, Vertex}

  \msgs{1.7}{5.4}{-1.1}{pergunta + \cd{session\_id}}
  \msgs{5.4}{9.1}{-1.85}{\cd{ask(message, session\_id)}}
  \msgself{9.1}{-2.5}{recupera memória da sessão (esquema \cd{adk})}
  \msgs{9.1}{12.8}{-3.4}{decide e invoca ferramenta}
  \msgs{12.8}{16.5}{-4.05}{SQL ou consulta vetorial}
  \msgr{16.5}{12.8}{-4.7}{dados estruturados}
  \msgself{12.8}{-5.35}{registra citações no coletor de contexto}
  \msgr{12.8}{9.1}{-6.2}{resultado JSON}
  \msgself{9.1}{-6.85}{repete se necessário (múltiplas ferramentas)}
  \msgs{9.1}{16.5}{-7.7}{geração final em contexto fechado}
  \msgr{16.5}{5.4}{-8.35}{texto em Markdown}
  \msgr{5.4}{1.7}{-9.1}{resposta + citações do coletor + selo de modo}

\end{tikzpicture}}
\caption[Sequência da conversa com agente]{Diagrama de sequência da conversa
  com agente. As citações devolvidas ao usuário vêm do coletor alimentado pelas
  ferramentas --- não do texto gerado na última etapa.}
\label{fig:seqchat}
\end{figure}
